\documentclass[openany]{ctexbook}
\usepackage{amsmath,amssymb}
\usepackage[hidelinks]{hyperref}
\newcommand{\note}[1]{\footnote{译注：#1}}
\title{代数几何基础第五卷预备笔记\\第8节末：域论与维数讨论}
\author{据用户上传英译编辑稿译出}
\date{2026年11月15日}
\begin{document}
\maketitle
\chapter*{范围与校订状态}
本批连续翻译底本第35页“From the [(cunutesque?) Fr] point of view and in terms of field theory”起，至第36页第9节标题前止，包括跨页域论段落及随后维数讨论、编排备忘。下一处为第36页第9节标题及9.1。本批为独立增补，未并入主合订本；不是对第8节全部数学结论的重新证明。8.8缺失条件仍未恢复，前批对偶记号及逆命题疑点仍然保留。
\chapter*{第8节末未编号讨论（续）}
从〔底本夹注：“(cunutesque?) Fr”，读法待核〕的观点出发，用域论的语言，可以如下用超越次数来表述这一条件。选择一个既不包含 $L^0$、也不包含 $X$ 的“无穷远超平面”，然后在其补集中讨论，也就是说，实质上在一个仿射型概形上讨论。\note{这里“包含 $X$”沿用底本措辞；在给定态射 $f:X\to P$ 的设置下，应按其像 $f(X)$ 理解，而不是假定 $f$ 已是嵌入。}
取在 $L^0$ 上消失的线性形式空间的一组基，记为 $T_1,\ldots,T_p$，其中 $p=\operatorname{codim}(L^0,P)$；考虑它们在 $X$ 的分式域 $K$ 中的拉回 $t_1,\ldots,t_p$（这里假定 $X$ 整）。\note{底本在这句写“vanishing on $L$”，但本段中心为 $L^0$，且紧接着指定 $p=\operatorname{codim}(L^0,P)$；这里统一中心记号为 $L^0$，区别于前段的超平面参数线性系 $L$。在所选仿射开集上，线性形式经平凡化成为函数，再拉回到函数域。}
至少有一个 $t_i$ 不为零，不妨设为 $t_1$。于是令
\[
a_1=t_2/t_1,\quad\ldots,\quad a_{p-1}=t_p/t_1.
\]
那么 $\dim g(X')$ 正是子域 $k(a_1,\ldots,a_{p-1})\subset K$ 对 $k$ 的超越次数。\note{底本写为“$K(a_1,\ldots,a_{p-1})\subset K$”。由于 $a_i\in K$，该域就是 $K$，一般不能给出像的维数。正文订为 $k(a_1,\ldots,a_{p-1})$：在投影的首坐标非零的仿射图中，像的坐标环由这些坐标比生成，其分式域即此子域，故像闭包的维数等于它对 $k$ 的超越次数。这一订正仅针对本句，并不恢复前批关于 $Q$ 的全部对偶约定。}
因此，若该超越次数至少为2，便符合要求。
若它等于1，则原稿要求：在 $k$ 上，$f(X)$ 包含于 $P$ 的某个包含 $L^0$、且维数至多为1的线性子簇中，并且 $g:X'\to g(X')$ 的泛纤维应当几何不可约。\note{“dimension at most one”是底本跨第35、36页的现存文字，不能直接当成已核实判据。若 $\dim L^0>1$，这样的线性子簇根本不能包含 $L^0$；而从中心 $L^0$ 投影后像落在直线的条件，通常对应于其逆像落在更高维的线性子空间。这里只明确引述原稿要求，未猜补可能漏失的相对维数或改写为未经证实的充要条件。}

设 $L^0$ 的维数为 $q$，则 $u:P-L^0\to Q$ 的纤维维数为 $q+1$。\note{底本先写“dimension $g$”，下一行及后续却均用 $q$；此处将维数变量统一为 $q$，不与态射 $g$ 混同。这里采用前文由中心 $L^0$ 给出的线性投影。}
将该投影限制到 $f(X)\setminus L^0$，其纤维维数至多为 $q+1$，因而
\[
\dim g(X')\ge \dim f(X)-(q+1).
\]
\note{底本写“so that those of $g$ are $\dim\le q+1$”。一般 $g=u\circ f'$ 的纤维还包含 $f$ 自身的纤维，不能仅由 $u$ 的纤维界直接得到这一说法。正文改用投影在像 $f(X)\setminus L^0$ 上的限制：其纤维是 $u$ 的纤维中的子集，维数界及随后关于像维数的不等式成立。若另知 $f$ 拟有限，原稿对 $g$ 的纤维界亦可成立；此处不擅自新增该假设。像的维数按其闭包的维数理解，并沿用 $X'=f^{-1}(P-L^0)$ 稠密的前提。}
所以，若 $\dim f(X)\ge q+3$，便满足关于 $g(X')$ 的维数条件。当 $q=0$ 时，就得到8.17(a)中所指出的事实。
回到8.18的条件，原稿指出：对于一个 $x\in T$，只要 $x$ 不是“相对于 $f$ 在 $x$ 处为锥形的”，条件(a)便满足，这里的意义是显然的。\note{底本为“provided $x$ is not ‘conical at $x$ relative to $f$’”。首个小写 $x$ 可能本应为 $X$，也可能是在描述点 $x$ 为相对于 $f$ 的锥点；对象和定义均未明，正文保留现存措辞，不确定恢复。原稿未在本段给出精确定义或完整论证，故这句不作为已经独立验证的判据。}
也许最好把这些最后的 Bertini 式讨论放到下一节中……更换射影浸入。
\end{document}
